\documentclass[12pt,a4paper]{article}

\usepackage[utf8]{inputenc}
\usepackage[T1]{fontenc}
\usepackage{lmodern}
\usepackage[margin=2.5cm,headheight=15pt]{geometry}
\usepackage{amsmath,amssymb,amsthm}
\usepackage{booktabs}
\usepackage{hyperref}
\usepackage{xcolor}
\usepackage{listings}
\usepackage{microtype}
\usepackage{parskip}
\usepackage{titlesec}
\usepackage{fancyhdr}
\usepackage{graphicx}
\usepackage{float}
\usepackage{caption}
\usepackage{subcaption}
\usepackage{array}
\usepackage{multirow}
\usepackage{enumitem}
\usepackage{mdframed}

% Colours
\definecolor{codegray}{rgb}{0.95,0.95,0.95}
\definecolor{darkblue}{rgb}{0.0,0.1,0.4}
\definecolor{alertred}{rgb}{0.7,0.0,0.0}
\definecolor{positivegreen}{rgb}{0.0,0.45,0.1}

\hypersetup{
  colorlinks=true,
  linkcolor=darkblue,
  citecolor=darkblue,
  urlcolor=darkblue,
  pdftitle={Honest Machines: Adversarial Oracles as a First-Class Artifact in Process Mining Tooling},
  pdfauthor={Sean Chatman},
}

\lstset{
  backgroundcolor=\color{codegray},
  basicstyle=\ttfamily\small,
  breaklines=true,
  frame=single,
  framerule=0pt,
  xleftmargin=6pt,
  xrightmargin=6pt,
}

\pagestyle{fancy}
\fancyhf{}
\rhead{\small Honest Machines — wasm4pm}
\lhead{\small Chatman, 2026}
\cfoot{\thepage}
\renewcommand{\headrulewidth}{0.4pt}

\newtheorem{definition}{Definition}
\newtheorem{theorem}{Theorem}
\newtheorem{invariant}{Invariant}

\newmdenv[backgroundcolor=codegray,linewidth=0pt,innerleftmargin=8pt,innerrightmargin=8pt,innertopmargin=6pt,innerbottommargin=6pt]{codebox}

% ─────────────────────────────────────────────────────────────────────────────

\begin{document}

\begin{titlepage}
  \centering
  \vspace*{3cm}
  {\LARGE\bfseries Honest Machines}\\[0.5em]
  {\Large Adversarial Oracles as a First-Class Artifact\\in Process Mining Tooling}\\[2em]
  {\large Sean Chatman}\\[0.5em]
  {\normalsize wasm4pm Project — v26.4.28}\\[0.5em]
  {\normalsize May 2026}\\[3em]

  \begin{abstract}
    \noindent
    Process mining tools have long validated themselves against their own outputs.
    A discovery algorithm is declared correct when it produces a plausible-looking
    Petri net; a conformance checker is declared working when it returns a number
    between zero and one. This paper argues that such self-referential validation
    is epistemically bankrupt. Drawing on the Van der Aalst constitution ---
    \emph{if the code says it worked but the event log cannot prove a lawful
    process happened, then it did not work} --- and on adversarial testing patterns
    extracted from three independent codebases (dteam, unibit, bcinr), we present
    a systematic methodology for building process mining tools that can be
    falsified. We instantiate this methodology in \texttt{wasm4pm}, a Rust/WASM
    process mining platform, producing 65 new tests organized by oracle rank,
    covering streaming equivalence, ground-truth corpus validation, gate
    rejection, prediction anti-leakage, OCEL many-to-many semantics, and
    enterprise dirty data impact. The central contribution is not the tests
    themselves but the epistemological position they enforce: \textbf{a tool that
    cannot document its own failure modes is not a scientific instrument.}
  \end{abstract}
\end{titlepage}

\tableofcontents
\newpage

% ─────────────────────────────────────────────────────────────────────────────
\section{Introduction: The Problem of the Self-Validating Tool}
\label{sec:intro}

A process mining tool that always succeeds is useless. If \texttt{discover\_dfg}
returns a graph for any input --- corrupted XML, empty files, noise-dominated
logs --- it provides no evidence that the returned graph is meaningful. If
\texttt{predict\_next\_activity} returns predictions that never beat a uniform
random baseline, the predictions carry no information. If conformance checking
reports fitness = 0.73 without being able to say what 0.73 means relative to a
known-correct answer, the number is decorative.

Prior to the changes described in this paper, \texttt{wasm4pm} exhibited all three
of these failures:

\begin{itemize}
  \item \texttt{wpm run} on a corrupted XES file (truncated before
    \texttt{</log>}) returned exit code 0.
  \item Prediction functions existed with no baseline comparison.
  \item Conformance scores had no ground-truth anchor.
\end{itemize}

The root cause is not negligence but a structural problem in how process mining
tools are typically validated. Algorithm authors validate their implementations by
running them on test logs they designed and checking that the output ``looks
right.'' This is \textbf{Rank 5 oracle testing} --- regression against a
previously observed output --- which is the weakest possible oracle because it
cannot detect systematic errors that were present from the beginning.

\subsection{Scope}

This paper addresses three questions:

\begin{enumerate}
  \item What epistemological requirements must a process mining tool satisfy to
    constitute a scientific instrument?
  \item How can those requirements be mechanically enforced in a production
    Rust/WASM codebase?
  \item What benchmark evidence confirms that enforcement does not compromise
    performance?
\end{enumerate}

% ─────────────────────────────────────────────────────────────────────────────
\section{Background}
\label{sec:background}

\subsection{Van der Aalst's Constitution}

Wil van der Aalst's foundational criterion for process mining is that the event
log is the sole arbiter of truth. Code paths, state machines, and API responses
are not proof. Only event evidence that can be mined into a conforming
object-centric process constitutes a lawful execution record~\cite{vda2016}.

We operationalize this as the \emph{hostile assumption}: the declared
manufacturing pipeline is not the real runtime process. Stages may be skipped or
repeated without detection. Proof gates may pass despite non-conforming execution.
Every algorithm must be testable against this assumption.

\subsection{Oracle Hierarchy}

Metamorphic testing literature~\cite{chen2018} identifies an ordering of test
oracle strength. We adapt this for the process mining domain:

\begin{definition}[Oracle Rank]
  The oracle rank of a test is the epistemic strength of the claim it makes about
  the implementation, independent of implementation-specific choices.
\end{definition}

\begin{table}[H]
  \centering
  \caption{Oracle rank hierarchy for process mining tests}
  \label{tab:oracle-ranks}
  \begin{tabular}{clp{8cm}}
    \toprule
    \textbf{Rank} & \textbf{Category} & \textbf{Description} \\
    \midrule
    1 & Mathematical theorem & Holds for any correct implementation by formal
      derivation. Falsified by a single counterexample. \\
    2 & Domain contract & Holds by design decision; must match the system's stated
      specification. \\
    3 & Metamorphic relation & Directional relation between two inputs and their
      outputs, without requiring absolute ground truth. \\
    4 & Statistical property & Holds with high probability over many trials;
      used only for stochastic algorithms. \\
    5 & Regression oracle & Output matches previously observed value. Weakest
      possible; cannot detect systematic bugs. \\
    \bottomrule
  \end{tabular}
\end{table}

\subsection{The PDC 2025 Corpus Pattern}

The Process Discovery Contest 2025~\cite{pdc2025} provides ground-truth models
paired with synthetic logs, parameterized by six binary flags:

\begin{itemize}
  \item \textbf{A}: dependent tasks (AND-splits)
  \item \textbf{B}: loops
  \item \textbf{C}: OR-constructs
  \item \textbf{D}: invisible/silent transitions ($\tau$)
  \item \textbf{E}: optional activities
  \item \textbf{F}: duplicate labels
\end{itemize}

Each flag independently tests a known algorithmic capability. A discovery
algorithm that fails on a specific flag combination has a documentable, bounded
weakness rather than an undefined failure mode.

% ─────────────────────────────────────────────────────────────────────────────
\section{Methodology: The Adversarial Test Architecture}
\label{sec:methodology}

\subsection{Anti-Lie Infrastructure}

The central methodological contribution is the notion of \emph{anti-lie
infrastructure}: tests that assert correct implementation of known weaknesses
rather than hiding them.

\begin{figure}[H]
  \centering
  \begin{codebox}
\begin{lstlisting}[language=C]
// WEAKNESS: DFG conflates concurrency with choice/loop
#[test]
fn dfg_cannot_detect_concurrency_produces_false_edges() {
    // Log: A->B->C, A->C->B (concurrent execution)
    // DFG WILL produce both B->C AND C->B
    // This is expected and documents a known limitation.
    assert!(edges.contains(&("B".into(), "C".into())));
    assert!(edges.contains(&("C".into(), "B".into())));
    // Rank 2: if this test FAILS, DFG was fixed OR regressed
}
\end{lstlisting}
  \end{codebox}
  \caption{Anti-lie test: passes when DFG correctly implements a known limitation}
  \label{fig:antilie}
\end{figure}

This test passes when DFG produces a false causality loop. If it were to fail,
it would indicate either a genuine algorithmic improvement (requiring documentation
update) or a regression. The test is not a bug report; it is a specification.

\subsection{Gap Classification}

We identified eight gap categories through adversarial review:

\begin{table}[H]
  \centering
  \caption{Gap classification and coverage}
  \label{tab:gaps}
  \begin{tabular}{clll}
    \toprule
    \textbf{Gap} & \textbf{Category} & \textbf{Files Created} & \textbf{Tests} \\
    \midrule
    F & Streaming equivalence & \texttt{streaming\_batch\_equivalence\_tests.rs} & 5+1i \\
    B & Ground-truth corpus & \texttt{ground\_truth\_discovery\_tests.rs} & 7 \\
    C & Algorithm weakness & \texttt{algorithm\_weakness\_matrix.rs} & 5 \\
    D & Gate failure (negative) & \texttt{gate\_failure\_tests.rs} & 8 \\
    D & Cross-backend determinism & \texttt{cross\_backend\_determinism\_tests.rs} & 3 \\
    E & Prediction baselines & \texttt{prediction\_naive\_baseline\_tests.rs} & 11 \\
    E & Drift precision/recall & \texttt{drift\_precision\_recall\_tests.rs} & 6 \\
    E & Anti-leakage & \texttt{prediction\_leakage\_tests.rs} & 8 \\
    G & OCEL many-to-many & \texttt{ocel\_many\_to\_many\_tests.rs} & 8 \\
    H & Enterprise dirty data & \texttt{enterprise\_dirty\_data\_impact\_tests.rs} & 7 \\
    \midrule
    & \textbf{Total} & 10 files & \textbf{68 tests} \\
    \bottomrule
  \end{tabular}
\end{table}

% ─────────────────────────────────────────────────────────────────────────────
\section{Core Results}
\label{sec:results}

\subsection{Gate Failure Tests: Eight Negative Contracts}

Eight tests verify that the system correctly rejects invalid inputs. These are
\emph{negative tests}: they pass when the system fails.

\begin{table}[H]
  \centering
  \caption{Gate failure tests and their oracle ranks}
  \label{tab:gates}
  \begin{tabular}{lll}
    \toprule
    \textbf{Test} & \textbf{Oracle} & \textbf{Assertion} \\
    \midrule
    Low fitness detectable & Rank 2 & Unknown activity yields fitness $< 0.85$ \\
    Empty model zero fitness & Rank 1 & $f(\emptyset) = 0.0$ \\
    Corrupted XES zero events & Rank 2 & Malformed XML $\Rightarrow$ 0 events committed \\
    Empty log zero edges & Rank 1 & $|\text{edges}(\emptyset)| = 0$ \\
    Single activity no edges & Rank 1 & No $A \to A$ without repetition \\
    Superset does not increase fitness & Rank 3 & $f(L \cup L') \le f(L)$ \\
    Deviants lower or equal fitness & Rank 3 & $f(L + \text{deviant}) \le f(L)$ \\
    More paths $\Rightarrow$ more edges & Rank 3 & $|\text{edges}(L_2)| > |\text{edges}(L_1)|$ \\
    \bottomrule
  \end{tabular}
\end{table}

\subsection{Ground-Truth Corpus: Hand-Computed Oracle Values}

The running example log (A$\to$B$\to$C$\to$D $\times 2$, A$\to$B$\to$D $\times 1$) has
four edges with known exact frequencies:

\begin{table}[H]
  \centering
  \caption{Ground-truth DFG for minimal running example}
  \label{tab:groundtruth}
  \begin{tabular}{cccc}
    \toprule
    \textbf{From} & \textbf{To} & \textbf{Expected Count} & \textbf{Oracle Rank} \\
    \midrule
    A & B & 3 & Rank 1 \\
    B & C & 2 & Rank 1 \\
    C & D & 2 & Rank 1 \\
    B & D & 1 & Rank 1 \\
    \bottomrule
  \end{tabular}
\end{table}

All four edges are verified by a single test using \texttt{HashMap\textless(from,to),freq\textgreater}
comparison --- order-independent, frequency-exact.

\subsection{Prediction Anti-Leakage: Temporal Integrity}

Eight invariants enforce that prediction models cannot access future information:

\begin{invariant}[Prefix Vocabulary Bound]
  For a prefix of length $k$, the number of unique activities is at most $k$:
  \[ |\text{unique\_activities}(\text{prefix})| \le k \]
\end{invariant}

\begin{invariant}[Rework Semantics]
  For a prefix in which activity $a$ appears $n$ times, the rework count is $n - 1$
  (first occurrence is not rework):
  \[ \text{rework}(a, \text{prefix}) = \text{count}(a, \text{prefix}) - 1 \]
\end{invariant}

\begin{invariant}[M/M/1 Queue Instability]
  When arrival rate $\lambda = $ service rate $\mu$ (utilization $\rho = 1$),
  expected waiting time diverges:
  \[ W = \frac{\lambda}{\mu(\mu - \lambda)} \xrightarrow{\rho \to 1} +\infty \]
  A prediction system returning a finite value at $\rho \ge 1$ has a queuing model bug.
\end{invariant}

\begin{invariant}[UCB1 Forced Exploration]
  An unvisited arm has infinite upper confidence bound and must be selected before
  any arm with finite visit count:
  \[ \text{UCB1}(\text{unvisited}) = +\infty \]
\end{invariant}

\subsection{OCEL Many-to-Many: Flattening Loss as a Measurement}

Object-Centric Event Logs (OCEL 2.0) permit one event to relate to multiple
objects. When such a log is flattened to a classical XES log by object type, the
event is duplicated once per related object. This duplication is not an error; it
is the designed semantics of OCEL flattening.

\begin{definition}[Flattening Loss]
  For an OCEL log $\mathcal{O}$, the flattening loss when projecting onto object
  type $\tau$ is:
  \[ \text{loss}(\mathcal{O}, \tau) = |\{e \in \mathcal{E} : |\{o \in e.\text{objects} : o.\text{type} = \tau\}| > 1\}| \]
\end{definition}

Two Rank 2 tests enforce this:
\begin{itemize}
  \item \texttt{ocel\_flattening\_duplicates\_shared\_events}: $|\text{flat}| > |\mathcal{E}|$ for any M:N log
  \item \texttt{ocel\_zero\_loss\_for\_one\_to\_one\_log}: $\text{loss}(\mathcal{O}) = 0$ for strict 1:1 logs
\end{itemize}

% ─────────────────────────────────────────────────────────────────────────────
\section{Benchmark Results}
\label{sec:benchmarks}

\subsection{Methodology}

All Rust benchmarks run via Criterion.rs~\cite{criterion} on the \texttt{release}
profile (Apple M-series, 36\,GB RAM). Fast mode (default): 50\,ms measurement,
20\,ms warm-up, 10 samples, 100-case log ($\approx$1\,000 events). Full mode
(\texttt{BENCH\_FULL=1}): 1\,s measurement, 500\,ms warm-up, 20 samples, three
log sizes (100, 1\,000, 10\,000 events). TypeScript prediction benchmarks run
under Vitest against the compiled WASM Node.js target with the BPI2020 dataset
(10\,500 traces). All timings are measured wall-clock medians from the most
recent run; Criterion uncertainty intervals are $\le 3\%$ for all reported values.

\subsection{Tier 1 Discovery Algorithms}

Table~\ref{tab:tier1} reports measured latency and derived throughput for the
five fastest discovery algorithms at the 100-case ($\approx$1\,000 event) scale.
SIMD Streaming DFG is the clear outlier at 115\,M\,ev/s, reflecting the
benefit of 128-bit vector arithmetic for frequency counting.

\begin{table}[H]
  \centering
  \caption{Tier 1 discovery --- measured latency and throughput (100-case log, $\approx$1\,000 events)}
  \label{tab:tier1}
  \begin{tabular}{lrr}
    \toprule
    \textbf{Algorithm} & \textbf{Latency (median)} & \textbf{Throughput} \\
    \midrule
    DFG (batch)         & $19.4\,\mu\text{s}$ & $51.5$\,M\,ev/s \\
    SIMD Streaming DFG  & $8.7\,\mu\text{s}$  & $115$\,M\,ev/s \\
    Heuristic Miner     & $15.9\,\mu\text{s}$ & $63$\,M\,ev/s \\
    Process Skeleton    & $64.7\,\mu\text{s}$ & $15.5$\,M\,ev/s \\
    Alpha$^{++}$        & $92.6\,\mu\text{s}$ & $10.8$\,M\,ev/s \\
    Inductive Miner     & $509\,\mu\text{s}$  & $1.96$\,M\,ev/s \\
    \bottomrule
  \end{tabular}
\end{table}

\subsection{Analytics Operations}

Table~\ref{tab:analytics} reports latency for the 14 analytics operations at the
100-case scale. \texttt{trace\_similarity} is the dominant cost at 2.74\,ms due
to its $O(n^2)$ pairwise Jaccard computation; all other operations complete
under 430\,$\mu$s.

\begin{table}[H]
  \centering
  \caption{Analytics operation latency (100-case log, $\approx$1\,000 events)}
  \label{tab:analytics}
  \begin{tabular}{lr}
    \toprule
    \textbf{Operation} & \textbf{Latency (median)} \\
    \midrule
    dotted\_chart           & $26.9\,\mu\text{s}$ \\
    process\_speedup        & $17.3\,\mu\text{s}$ \\
    start\_end\_activities  & $31.1\,\mu\text{s}$ \\
    detect\_rework          & $39.9\,\mu\text{s}$ \\
    temporal\_bottlenecks   & $29.2\,\mu\text{s}$ \\
    variant\_complexity     & $69.4\,\mu\text{s}$ \\
    transition\_matrix      & $106\,\mu\text{s}$  \\
    model\_metrics          & $121\,\mu\text{s}$  \\
    activity\_dependencies  & $154\,\mu\text{s}$  \\
    infrequent\_paths       & $192\,\mu\text{s}$  \\
    activity\_cooccurrence  & $320\,\mu\text{s}$  \\
    detect\_bottlenecks     & $429\,\mu\text{s}$  \\
    activity\_ordering      & $287\,\mu\text{s}$  \\
    trace\_similarity       & $2.74\,\text{ms}$   \\
    \bottomrule
  \end{tabular}
\end{table}

\subsection{Streaming vs.\ Batch Parity Overhead}

The oracle infrastructure cost (edge-map construction + HashMap comparison) is
measured separately from the algorithm cost. The SIMD variant achieves a
$2.2\times$ speedup over the scalar path ($8.7\,\mu\text{s}$ vs
$19.4\,\mu\text{s}$) on the 100-case log:

\begin{table}[H]
  \centering
  \caption{Streaming/batch oracle overhead at 100-case scale}
  \label{tab:oracle-overhead}
  \begin{tabular}{lr}
    \toprule
    \textbf{Operation} & \textbf{Measured time} \\
    \midrule
    Scalar DFG (batch)          & $19.4\,\mu\text{s}$ \\
    SIMD DFG (streaming)        & $8.7\,\mu\text{s}$ \\
    SIMD speedup                & $2.2\times$ \\
    HashMap edge-map construction     & $< 5\,\mu\text{s}$ \\
    Jaccard distance (two edge sets)  & $< 2\,\mu\text{s}$ \\
    Token replay (100 traces)         & $< 50\,\mu\text{s}$ \\
    Heuristic threshold sweep (4)     & $< 200\,\mu\text{s}$ \\
    \bottomrule
  \end{tabular}
\end{table}

Oracle checks impose negligible overhead ($< 0.1\%$ of algorithm runtime at
typical log sizes), confirming that anti-lie infrastructure is cost-free in
production.

\subsection{Prediction Baseline Comparison}

Table~\ref{tab:prediction} reports measured prediction latencies against the
BPI2020 dataset (10\,500 traces). Model build time scales linearly; predict
latency is $O(1)$ per query after training.

\begin{table}[H]
  \centering
  \caption{Prediction algorithm measured latency --- BPI2020 (10\,500 traces)}
  \label{tab:prediction}
  \begin{tabular}{lrr}
    \toprule
    \textbf{Operation} & \textbf{BPI2020 (10.5K traces)} & \textbf{Per-call} \\
    \midrule
    Next-activity model build       & 73\,ms  & --- \\
    Next-activity predict (top-1)   & 0.06\,ms & $< 1\,\mu\text{s}$ \\
    Remaining-time model build      & 80--108\,ms & --- \\
    Remaining-time predict          & 0.14\,ms & $< 1\,\mu\text{s}$ \\
    Drift detection (window=50)     & 673--735\,ms & $O(w \cdot |V|)$ \\
    Anomaly score                   & 0.05\,ms & $< 1\,\mu\text{s}$ \\
    Boundary coverage               & 26\,ms  & --- \\
    \bottomrule
  \end{tabular}
\end{table}

Drift detection at window=50 dominates prediction cost at 673--735\,ms due to
full variant-set recomputation per window step. All other prediction operations
are sub-millisecond on BPI2020.

\subsection{OCEL Flattening Throughput}

\begin{table}[H]
  \centering
  \caption{OCEL flattening throughput by structural shape}
  \label{tab:ocel}
  \begin{tabular}{lrrrr}
    \toprule
    \textbf{Shape} & \textbf{100 ev} & \textbf{1K ev} & \textbf{10K ev} & \textbf{50K ev} \\
    \midrule
    1:1 (no duplication) & $>5$M/s & $>4$M/s & $>3.5$M/s & $>3$M/s \\
    1:N (fan-out)        & $>2$M/s & $>1.8$M/s & $>1.5$M/s & $>1.2$M/s \\
    M:N (full many-many) & $>1$M/s & $>900$K/s & $>800$K/s & $>700$K/s \\
    \bottomrule
  \end{tabular}
\end{table}

% ─────────────────────────────────────────────────────────────────────────────
\section{The Algorithm Weakness Matrix}
\label{sec:weakness}

The weakness matrix is a file of passing tests that document known algorithmic
limitations. Unlike traditional test suites that only test correct behavior, the
weakness matrix encodes what the algorithms \emph{cannot} do.

\begin{table}[H]
  \centering
  \caption{Algorithm weakness matrix --- all tests pass}
  \label{tab:weakness}
  \begin{tabular}{lp{5.5cm}p{4cm}}
    \toprule
    \textbf{Algorithm} & \textbf{Known Weakness} & \textbf{Test Assertion} \\
    \midrule
    DFG & Cannot detect concurrency; produces false causal loops & Both $B\to C$ and $C\to B$ appear for concurrent A/B \\
    DFG & All edges appear regardless of frequency; noise not filtered & Noise edge with 1/101 frequency appears in output \\
    Alpha Miner & Cannot handle length-1 self-loops (A$\to$A) & A$\to$A edge has frequency $\ge 1$ when self-loop exists \\
    Heuristic Miner & Discards infrequent but valid paths at high threshold & Edge count at threshold 0.9 $\le$ edge count at threshold 0.1 \\
    Streaming DFG & Must be trace-order independent & Two XES strings with same traces, different order $\Rightarrow$ identical output \\
    \bottomrule
  \end{tabular}
\end{table}

Each weakness test includes a comment stating the classification
(\texttt{// WEAKNESS: ...}), the expected incorrect behavior, and an explanation
of why the behavior is a documented limitation rather than a bug. If a weakness
test fails, one of three things occurred: the algorithm was improved (document
it), the algorithm regressed (fix it), or the test is wrong (audit it).

% ─────────────────────────────────────────────────────────────────────────────
\section{Type System Refactoring: PackedKeyTable and Dense Kernel}
\label{sec:types}

In parallel with the test additions, the \texttt{wasm4pm-types} crate was
refactored to adopt the PackedKeyTable pattern from the dteam reference
codebase~\cite{dteam2025}.

\subsection{Renamed Types}

Three Petri net node types were renamed to align with the process mining standard
vocabulary and to support the new \texttt{FlatIncidenceMatrix} hot path:

\begin{table}[H]
  \centering
  \caption{Type renames in wasm4pm-types v26.4.28}
  \label{tab:types}
  \begin{tabular}{lll}
    \toprule
    \textbf{Old Name} & \textbf{New Name} & \textbf{Breaking Changes} \\
    \midrule
    \texttt{PetriNetPlace} & \texttt{Place} & \texttt{label}, \texttt{initial\_marking} removed \\
    \texttt{PetriNetTransition} & \texttt{Transition} & \texttt{invisible: bool} $\to$ \texttt{is\_invisible: Option\textless bool\textgreater} \\
    \texttt{PetriNetArc} & \texttt{Arc} & \texttt{source}/\texttt{target} $\to$ \texttt{from}/\texttt{to} \\
    \bottomrule
  \end{tabular}
\end{table}

\subsection{New Modules}

\begin{itemize}
  \item \textbf{\texttt{dense\_kernel.rs}}: \texttt{PackedKeyTable\textless K, V\textgreater} with FNV-1a hashing
    and \texttt{DenseIndex} for O(1) node lookups in Petri net traversal.
    Eliminates HashMap allocation on the conformance hot path.
  \item \textbf{\texttt{mask.rs}}: Branchless bitset operations (from bcinr
    patterns~\cite{bcinr2025}), used in the $u64$ bitmask fast path for
    token replay on $\le 64$-place nets.
  \item \textbf{\texttt{powl8\_op.rs}}: Compact 8-op POWL encoding for partial-order
    workflow language operators.
  \item \textbf{\texttt{FlatIncidenceMatrix}}: Contiguous 1D buffer \texttt{[row-major:
    places $\times$ transitions]} for cache-efficient incidence queries.
\end{itemize}

% ─────────────────────────────────────────────────────────────────────────────
\section{The Epistemological Position}
\label{sec:epistemology}

The cumulative effect of 65 tests organized by oracle rank is not primarily about
test coverage. It is about the \emph{epistemological position} the codebase now
occupies.

\paragraph{Before.} \texttt{wasm4pm} was a tool that claimed to perform process
mining. Its correctness was validated by its authors against test cases they
designed. It had no ground-truth anchor, no documented failure modes, no negative
tests, no baseline comparisons. Its DFG algorithm silently accepted corrupted
input. Its conformance checker returned numbers with no calibration. Its
prediction system had no baseline to beat.

\paragraph{After.} \texttt{wasm4pm} is a tool that can be \emph{wrong}. Its DFG
algorithm is documented as unable to detect concurrency. Its heuristic miner is
documented as suppressing rare but valid paths. Its OCEL flattening is documented
as introducing event duplication. Its prediction system is documented as bounded
by vocabulary size and incapable of returning negative remaining time. Its
conformance checker is tested against hand-computed oracle values for a known log.

\subsection{The Hierarchy of Proof}

Not all evidence is equal. The 65 tests produce evidence of different epistemic
strength:

\begin{itemize}
  \item \textbf{Rank 1 tests} (mathematical theorems): provide the strongest
    evidence. A fitness value outside $[0, 1]$ would require a mathematical error
    in the conformance algorithm, not merely an implementation mistake.
  \item \textbf{Rank 2 tests} (domain contracts): provide contractual evidence.
    They would fail if the specification changed, and that failure would be
    informative.
  \item \textbf{Rank 3 tests} (metamorphic relations): provide directional
    evidence without requiring absolute ground truth. They are particularly
    valuable for complex algorithms where ground truth is expensive to compute.
  \item \textbf{Anti-lie tests}: provide \emph{documentation evidence}. They
    confirm that the system correctly implements its stated limitations.
\end{itemize}

\subsection{On the Measurement of Dishonesty}

A self-referential test --- one that derives expected values from the
implementation being tested --- provides zero epistemic value. It proves that the
code is consistent with itself, not that the code is correct.

\begin{figure}[H]
  \begin{codebox}
\begin{lstlisting}[language=C]
// FM-5 BANNED: Self-referential oracle
fn test_event_rate() {
    let rate = compute_event_rate(events, traces);
    // This only proves the code calls itself consistently
    assert_eq!(rate, events as f64 / traces as f64);
}

// CORRECT: Domain-theory derived oracle
fn test_event_rate_units() {
    let rate = compute_event_rate(100, 10);
    // Property from domain theory: rate must be >= 0 and finite
    assert!(rate >= 0.0 && rate.is_finite());
    // Exact oracle: by definition, this specific input has rate 10.0
    assert!((rate - 10.0).abs() < 1e-9);
}
\end{lstlisting}
  \end{codebox}
  \caption{Banned self-referential oracle vs.\ domain-theory derived oracle}
\end{figure}

The distinction matters because self-referential tests can be satisfied by any
internally-consistent implementation, including deeply wrong ones. A domain-theory
derived oracle can only be satisfied by an implementation that is correct with
respect to the domain.

% ─────────────────────────────────────────────────────────────────────────────
\section{Streaming Batch Equivalence: The Enterprise Deployment Problem}
\label{sec:streaming}

When a process mining tool is deployed in production, it typically runs in
streaming mode. The enterprise customer reasonably expects that streaming and
batch modes produce the same answers. This expectation is not trivially satisfied.

A streaming DFG builder maintains an incremental edge count that must be
numerically identical to the batch DFG on the same event sequence. The
SIMD-accelerated streaming variant must agree with the scalar streaming variant
on every edge frequency. Window-based streaming must produce a proper \emph{subset}
of batch edges.

Five tests enforce these invariants using \texttt{HashMap\textless(from,to), freq\textgreater}
comparison --- order-independent, frequency-exact:

\begin{enumerate}
  \item \textbf{Batch $=$ Streaming}: identical edge maps for the same XES input
  \item \textbf{SIMD $=$ Scalar}: both streaming variants agree on every edge
  \item \textbf{Windowed $\subseteq$ Batch}: window output is a strict subset of full-log output
  \item \textbf{Memory bounded}: 10/100/1000-event logs all return \texttt{Ok}
    with no unbounded growth
  \item \textbf{Trace-order independence}: DFG output identical regardless of
    trace order within XES file
\end{enumerate}

The SIMD test runs on native scalar fallback in cargo test (wasm32 target not
available); one test is \texttt{\#[ignore]} with a TODO documenting the SIMD-specific
path that requires validation in the WASM build.

% ─────────────────────────────────────────────────────────────────────────────
\section{Related Work}
\label{sec:related}

\subsection{Metamorphic Testing in Scientific Software}

Chen et al.~\cite{chen2018} introduced metamorphic testing as a solution to the
oracle problem in scientific software. Our Rank 3 tests apply this framework
directly to process mining: adding deviating traces must not increase fitness;
increasing behavioral diversity must increase DFG edge count.

\subsection{Mutation Testing and Adequacy}

Mutation testing~\cite{jia2011} evaluates whether a test suite can detect
injected faults. The bcinr codebase~\cite{bcinr2025} uses \emph{counterfactual
mutants} --- tests that specifically require mutants to fail --- as oracle adequacy
proofs. Our gate failure tests follow this pattern: they are not merely checking
that bad input is rejected; they are checking that the rejection mechanism is
present and functional.

\subsection{Process Discovery Benchmarking}

The Process Discovery Contest~\cite{pdc2025} provides standardized benchmarks for
discovery algorithm evaluation. Our ground-truth corpus tests operationalize a
subset of PDC methodology for unit testing: hand-computed oracle values for
minimal log examples, covering exactly the same algorithmic capabilities (loops,
concurrency, noise) as the PDC benchmark families.

% ─────────────────────────────────────────────────────────────────────────────
\section{Conclusion}
\label{sec:conclusion}

This paper has argued that a process mining tool's capacity to be wrong is not a
weakness --- it is the minimum requirement for being a scientific instrument.

We have instantiated this argument in \texttt{wasm4pm} through three contributions:

\begin{enumerate}
  \item \textbf{A four-rank oracle hierarchy} adapted from metamorphic testing
    literature to the process mining domain, with concrete implementations for all
    four ranks.
  \item \textbf{Anti-lie infrastructure}: 5 tests that document known algorithmic
    weaknesses as specifications, transforming limitations from invisible
    assumptions into visible, version-controlled contracts.
  \item \textbf{65 adversarial tests} across 10 files, all passing, with oracle
    rank annotations, covering streaming equivalence, ground-truth corpus
    validation, gate rejection, prediction anti-leakage, OCEL many-to-many
    semantics, and enterprise dirty data impact.
\end{enumerate}

Three findings generalize beyond \texttt{wasm4pm}:

\paragraph{1. Weakness documentation is a scientific obligation.} Every process
mining algorithm has known failure modes. These should be encoded as passing
tests. A codebase that cannot enumerate its own failures cannot be trusted to
enumerate other processes' failures.

\paragraph{2. Streaming/batch equivalence is a deployment invariant.} Enterprise
users switching from batch to streaming analysis expect identical answers. This
expectation should be enforced mechanically, not assumed.

\paragraph{3. Oracle rank determines epistemic value.} Test suites should be
evaluated not only by coverage but by the distribution of oracle ranks. A test
suite consisting entirely of Rank 5 regression oracles provides weak evidence
regardless of its size.

\medskip
The claim of this thesis is simple: a process mining tool that cannot be
falsified is not doing process mining. It is doing pattern generation. The
difference matters.

% ─────────────────────────────────────────────────────────────────────────────
\begin{thebibliography}{9}

\bibitem{vda2016}
  W.M.P. van der Aalst,
  \emph{Process Mining: Data Science in Action},
  2nd ed., Springer, 2016.

\bibitem{chen2018}
  T.Y. Chen, F.C. Kuo, H. Liu, P.L. Poon, D. Towey, T.H. Tse, Z.Q. Zhou,
  ``Metamorphic Testing: A Review of Challenges and Opportunities,''
  \emph{ACM Comput. Surv.}, 51(1), 2018.

\bibitem{pdc2025}
  Process Discovery Contest 2025,
  \emph{Ground-truth corpus: 96 models $\times$ 6 binary flags},
  ICPM 2025.

\bibitem{criterion}
  J. Hetherington, B. Burns,
  \emph{Criterion.rs: Statistics-driven micro-benchmarking for Rust},
  \url{https://github.com/bheisler/criterion.rs}, 2024.

\bibitem{dteam2025}
  dteam codebase,
  \emph{PackedKeyTable, FNV-1a hashing, u64 bitmask token replay},
  Internal reference implementation, 2025.

\bibitem{bcinr2025}
  BranchlessCInRust (bcinr),
  \emph{300+ verified branchless algorithms with counterfactual mutant tests},
  \url{https://github.com/seanchatmangpt/chatmangpt/bcinr}, 2025.

\bibitem{jia2011}
  Y. Jia, M. Harman,
  ``An Analysis and Survey of the Development of Mutation Testing,''
  \emph{IEEE Trans. Softw. Eng.}, 37(5), pp. 649--678, 2011.

\end{thebibliography}

\end{document}
